\lysh{22062_SOLUTION}{1}
ΛΥΣΗ

α) Γενικά, η εξίσωση $(x - x_0)^2 + (y - y_0)^2 = \rho^2$ παριστά κύκλο με κέντρο $K(x_0, y_0)$ και ακτίνα $\rho$, αν $\rho > 0$ (ή ακτίνα $|\rho|$, αν $\rho \ne 0$). Συνεπώς, αν $\lambda \ne 0$, η εξίσωση $(x - \lambda)^2 + (y - \lambda)^2 = \lambda^2$ παριστά κύκλο με κέντρο $K(\lambda, \lambda)$ και ακτίνα $|\lambda|$.

β) Παρατηρούμε ότι για $x = \lambda$ και $y = \lambda$ επαληθεύεται η εξίσωση $y = x$. Επομένως, για κάθε $\lambda \ne 0$, το κέντρο $K(\lambda, \lambda)$ του κύκλου $C_\lambda$ είναι σημείο της ευθείας $y = x$.

Για τα ερωτήματα γ), δ), ε) υπενθυμίζουμε ότι μία ευθεία $\varepsilon: Ax + By + \Gamma = 0$ (με $|A| + |B| \ne 0$) εφάπτεται σε ένα κύκλο με κέντρο $K(x_0, y_0)$ και ακτίνα $\rho$ αν και μόνο αν $d(K, \varepsilon) = \rho$ ή ισοδύναμα αν
\[
\frac{|Ax_0+By_0+\Gamma|}{\sqrt{A^2+B^2}} = \rho \quad (1).
\]
Επίσης, από το α) ερώτημα, για κάθε $\lambda \ne 0$, η εξίσωση $(x - \lambda)^2 + (y - \lambda)^2 = \lambda^2$ παριστά κύκλο $C_\lambda$ με κέντρο $K(\lambda, \lambda)$ και ακτίνα $|\lambda|$. Οπότε:
\[
x_0 = \lambda, \quad y_0 = \lambda, \quad \rho = |\lambda|
\]

γ) Θεωρούμε την ευθεία $x = 0$. Εδώ: $A = 1, B = 0, \Gamma = 0$. Οπότε, η σχέση (1) ικανοποιείται για κάθε $\lambda \ne 0$, αφού:
\[
\frac{|1 \cdot \lambda + 0 \cdot \lambda + 0|}{\sqrt{1^2 + 0^2}} = |\lambda| \quad (\text{ταυτότητα}).
\]
Άρα, η ευθεία $x = 0$ εφάπτεται σε όλους τους κύκλους $C_\lambda, \lambda \ne 0$.

Για να δείξουμε ότι η ευθεία $y = 0$ εφάπτεται σε όλους τους κύκλους $C_\lambda$ θα μπορούσαμε να εργαστούμε όμοια με παραπάνω ή να παρατηρήσουμε ότι υπάρχει συμμετρία ως προς την $1^\eta$ διχοτόμο, καθώς η εξίσωση $(x - \lambda)^2 + (y - \lambda)^2 = \lambda^2$ παραμένει αναλλοίωτη αν εναλλάξουμε τους ρόλους των $x$ και $y$.

δ) Έστω $\alpha \ne 0$. Θεωρούμε την ευθεία $x = \alpha$. Εδώ: $A = 1, B = 0, \Gamma = -\alpha$. Οπότε, η σχέση (1) γίνεται:
\[
\frac{|1 \cdot \lambda + 0 \cdot \lambda - \alpha|}{\sqrt{1^2 + 0^2}} = |\lambda| \Leftrightarrow |\lambda - \alpha| = |\lambda| \Leftrightarrow \lambda - \alpha = \pm \lambda.
\]
Το θετικό πρόσημο δίνει $\alpha = 0$ (πράγμα αδύνατο, αφού έχει υποτεθεί $\alpha \ne 0$), ενώ το αρνητικό πρόσημο δίνει την μοναδική αποδεκτή τιμή του $\lambda$ που είναι το $\frac{\alpha}{2}$.

Λόγω συμμετρίας, το ίδιο συμβαίνει και για την ευθεία $y = \alpha$.
\end{lysh}